% ============================================================
% Unidade 1 — Parte B: erros, observações, exemplos, questões, síntese, respostas
% Fonte única: dossies/u1.md · respostas conferidas em verif/u1.py
% ============================================================
\ErrosComuns

\begin{erro}{Esquecer o universo}{1-E1}
\campo{O erro.} Resolver $x^2=2$ e responder $\{-\sqrt2,\sqrt2\}$ quando o enunciado procurava entre as frações; ou responder $\{-2,3\}$ para $x^2-x-6=0$ quando o universo era $\mathbb{N}$, onde $-2$ não está.
\campo{Por que parece plausível.} A equação é a mesma; resolver \enquote{dá} as raízes, e o universo parece um detalhe do enunciado.
\campo{Onde quebra.} $\{x\in\mathbb{Q}\mid x^2=2\}=\varnothing$, mas $\{x\in\mathbb{R}\mid x^2=2\}=\{-\sqrt2,\sqrt2\}$. E $\{x\in\mathbb{Z}\mid 2x+1=0\}=\varnothing$, mas em $\mathbb{Q}$ o conjunto é $\{-\tfrac12\}$.
\campo{Como detectar.} Depois de escrever a resposta, confira cada elemento: ele pertence a $U$? E falta algum elemento de $U$ que satisfaz a condição?
\campo{Como corrigir.} Escreva o universo na primeira linha da resolução, $\{x\in U\mid\ldots\}$, e filtre as soluções no fim (construção 1-H3).
\campo{Onde aparece.} \exref{ex:u1:universo}; questões \ref{u1:unitarios} e \ref{u1:equacao}.
\end{erro}

\begin{erro}{Reticências sem lei clara}{1-E2}
\campo{O erro.} Escrever $\{1,2,4,\ldots\}$ ou $\{3,5,7,\ldots\}$ e supor que todo leitor vai continuar do mesmo jeito.
\campo{Por que parece plausível.} Quem escreveu já tinha uma lei na cabeça; para ele, \enquote{o padrão é óbvio}.
\campo{Onde quebra.} Em $\{1,2,4,\ldots\}$, dobrando, o próximo é 8; somando 1, depois 2, depois 3, o próximo é 7. Em $\{3,5,7,\ldots\}$, os ímpares a partir de 3 contêm o 9; os primos ímpares, não.
\campo{Como detectar.} Tente escrever duas leis diferentes que comecem com os mesmos termos. Se conseguir, a lista não fixou o conjunto.
\campo{Como corrigir.} Liste termos suficientes para eliminar as outras leis ou, melhor, troque a lista pela propriedade (construção 1-H2).
\campo{Onde aparece.} \exref{ex:u1:reticencias}; \qref{u1:reticencias}.
\end{erro}

\begin{erro}{Confundir $\varnothing$ com $\{0\}$ ou com $\{\varnothing\}$}{1-E3}
\campo{O erro.} Dizer que $\{0\}$ é vazio (\enquote{só tem zero}), que $\{\varnothing\}$ é vazio, ou que $\{x\mid x-x=0\}$ é vazio.
\campo{Por que parece plausível.} \enquote{Zero} e \enquote{vazio} lembram \enquote{nada}; e $x-x=0$ \enquote{dá zero}.
\campo{Onde quebra.} $\{0\}$ tem um elemento, o número 0; $\{\varnothing\}$ tem um elemento, o conjunto vazio. E $x-x=0$ é verdadeira para \emph{todo} $x$: $\{x\in U\mid x-x=0\}=U$. Vazio é $\{x\in U\mid x-x=1\}$.
\campo{Como detectar.} Conte os objetos entre as chaves. Para uma propriedade, pergunte: existe algum candidato que a satisfaz?
\campo{Como corrigir.} $\varnothing$ é o único conjunto sem elementos; qualquer coisa escrita entre as chaves já é um elemento (construção 1-H4).
\campo{Onde aparece.} \exref{ex:u1:vazio}; questões \ref{u1:quantos} e \ref{u1:vazios}.
\end{erro}

\begin{erro}{Aceitar uma descrição sem critério de pertinência}{1-E4}
\campo{O erro.} Tratar \enquote{os números grandes}, \enquote{os bons alunos} ou \enquote{as cidades perto do mar} como conjuntos.
\campo{Por que parece plausível.} Na conversa, todo mundo entende mais ou menos o que se quer dizer.
\campo{Onde quebra.} O número $10^6$ é grande? E $10^3$? Cada leitor decide de um jeito, e não há como responder \enquote{$x\in A$?} para todo $x$. Já \enquote{os naturais menores que $10^{100}$} é um conjunto: enorme, mas decidido.
\campo{Como detectar.} Escolha um candidato de fronteira e pergunte se ele pertence. Se a resposta depende de opinião, não há conjunto.
\campo{Como corrigir.} Troque a palavra vaga por uma condição decidível: \enquote{média maior ou igual a 7}, \enquote{a menos de 50 km do litoral} (construção 1-H1 e caixa \emph{Para lembrar}).
\campo{Onde aparece.} \exref{ex:u1:classificar}; \qref{u1:criterio}.
\end{erro}

% ------------------------------------------------------------
\Observacoes

\begin{obs}{1-O1}
\textbf{O conjunto como objeto.} Para os pioneiros, não era óbvio que uma coleção infinita, como $\mathbb{N}$, pudesse ser tratada como uma coisa só, sobre a qual se opera, do mesmo modo que o número 3. O historiador José Ferreirós aponta essa \enquote{objetidade} como o elemento novo da teoria dos conjuntos: as ideias de coleção e de classe são muito antigas; o que mudou foi tratar o conjunto como um objeto matemático.\src{sepEarly}
\end{obs}

\begin{obs}{1-O2}
\textbf{Não existe \enquote{o conjunto de tudo}.} O axioma de separação tem uma consequência curiosa. Dado qualquer conjunto $S$, o conjunto $\{x\in S\mid x\notin x\}$ não pode ser elemento de $S$ (é o argumento de Russell, agora sem contradição). Portanto, nenhum conjunto contém todos os objetos: Zermelo concluiu, em 1908, que o domínio de todos os objetos não é, ele próprio, um conjunto.\src{sepZermelo}\src{zermelo1908} É por isso que o universo $U$ de um problema é sempre \enquote{um universo}, escolhido para aquele assunto, e nunca \enquote{o universo de tudo}.
\end{obs}

% ------------------------------------------------------------
\Exemplos

\begin{exemplo}{}{1}{adaptado do FME, exercícios 11 e 12}{1-H1, 1-H2, 1-U1}
\label{ex:u1:listar}
\begin{enumerate}[label=\alph*)]
  \item Dê os elementos de $A=\{x\mid x\text{ é letra da palavra \emph{matemática}}\}$ e de $B=\{x\mid x\text{ é cor da bandeira brasileira}\}$.
  \item Descreva por uma propriedade característica $C=\{0,2,4,6,8,\ldots\}$ e $D=\{0,1,2,\ldots,9\}$.
\end{enumerate}
\plano Em (a), percorrer os candidatos (as letras da palavra, as cores da bandeira) e guardar cada um uma vez. Em (b), achar o que os listados têm em comum e mais ninguém tem.
\resolucao (a) $A=\{m,a,t,e,i,c\}$ e $B=\{\text{verde},\text{amarelo},\text{azul},\text{branco}\}$. (b) $C=\{x\in\mathbb{N}\mid x\text{ é par}\}$ e $D=\{x\in\mathbb{N}\mid x\le 9\}$ (os algarismos).
\comentario A palavra \emph{matemática} tem dez letras, mas $A$ tem seis elementos: a letra \emph{a} aparece três vezes na palavra, e \emph{m} e \emph{t} aparecem duas vezes cada; no conjunto, cada uma entra uma vez só, porque o conjunto registra apenas \emph{se} a letra pertence (a Unidade~2 volta a isso). Em (b), \enquote{par} sozinho não bastaria para $C$: os pares negativos ficariam dentro; o universo $\mathbb{N}$ faz o trabalho.
\end{exemplo}

\begin{exemplo}{}{1}{autoral}{1-H1, 1-H4, 1-E3}
\label{ex:u1:vazio}
Quantos elementos têm $\varnothing$, $\{0\}$, $\{\varnothing\}$ e $\{x\in\mathbb{R}\mid x\neq x\}$?
\plano Contar o que está entre as chaves; para a propriedade, procurar algum real que a satisfaça.
\resolucao $\varnothing$ não tem elementos: 0. $\{0\}$ tem um elemento, o número 0: 1. $\{\varnothing\}$ tem um elemento, o conjunto vazio: 1. Nenhum real é diferente de si mesmo, então $\{x\in\mathbb{R}\mid x\neq x\}=\varnothing$: 0.
\comentario Os quatro números são 0, 1, 1, 0. O erro 1-E3 está em achar que $\{0\}$ ou $\{\varnothing\}$ são vazios: os dois têm um elemento, e o fato de um conjunto poder ser elemento de outro (construção 1-H1) é o que torna $\{\varnothing\}$ diferente de $\varnothing$.
\end{exemplo}

\begin{exemplo}{}{2}{autoral}{1-H3, 1-U2, 1-E1}
\label{ex:u1:universo}
Determine $\{x\in U\mid x^2-2=0\}$ para $U=\mathbb{Q}$ e para $U=\mathbb{R}$; e $\{x\in U\mid 2x+1=0\}$ para $U=\mathbb{Z}$ e para $U=\mathbb{Q}$.
\plano Resolver cada equação uma vez, como se não houvesse universo, e depois guardar só as soluções que estão em $U$.
\resolucao As soluções de $x^2=2$ são $\sqrt2$ e $-\sqrt2$, que não são frações. Logo $\{x\in\mathbb{Q}\mid x^2-2=0\}=\varnothing$ e $\{x\in\mathbb{R}\mid x^2-2=0\}=\{-\sqrt2,\sqrt2\}$. A solução de $2x+1=0$ é $-\tfrac12$, que não é inteiro: $\{x\in\mathbb{Z}\mid 2x+1=0\}=\varnothing$ e $\{x\in\mathbb{Q}\mid 2x+1=0\}=\{-\tfrac12\}$.
\comentario A mesma condição deu conjuntos diferentes. Por isso a construção 1-H3 manda fixar $U$ antes de resolver, e o teste do erro 1-E1 é conferir, no fim, se cada elemento pertence a $U$.
\end{exemplo}

\begin{exemplo}{}{2}{autoral}{1-H2, 1-U1, 1-E2}
\label{ex:u1:reticencias}
Encontre duas leis diferentes para a lista $\{1,2,4,\ldots\}$, escreva cada conjunto por uma propriedade e diga qual número viria depois do 4 em cada caso.
\plano Procurar leis que produzam 1, 2, 4 e depois se separem; escrever cada uma na forma \enquote{expressão $\mid$ condição}.
\resolucao Primeira lei: cada termo é o dobro do anterior; o conjunto é $\{2^n\mid n\in\mathbb{N}\}=\{1,2,4,8,16,\ldots\}$, e depois do 4 vem o 8. Segunda lei: soma-se 1, depois 2, depois 3, e assim por diante; o termo de ordem $n$ (contando de $n=0$) é $1+\tfrac{n(n+1)}{2}$, e o conjunto é $\{1+\tfrac{n(n+1)}{2}\mid n\in\mathbb{N}\}=\{1,2,4,7,11,\ldots\}$, em que depois do 4 vem o 7.
\comentario Os dois conjuntos começam iguais e são diferentes: o 7 está em um e não no outro. A lista $\{1,2,4,\ldots\}$ não fixou nenhum dos dois (erro 1-E2). A propriedade fixou.
\end{exemplo}

\begin{exemplo}{}{2}{autoral}{1-H4, 1-U3, 1-E4}
\label{ex:u1:classificar}
Classifique como vazio, unitário, finito com mais de um elemento ou infinito, ou diga que a descrição não define conjunto:
\begin{multicols}{2}
\begin{enumerate}[label=\alph*)]
  \item $\{x\in\mathbb{N}\mid x<0\}$
  \item $\{x\in\mathbb{Z}\mid x^2=9\text{ e }x>0\}$
  \item $\{x\in\mathbb{Z}\mid 3<x<4\}$
  \item $\{x\in\mathbb{R}\mid 3<x<4\}$
  \item os números grandes
\end{enumerate}
\end{multicols}
\plano Para cada item, ver quantos candidatos do universo satisfazem a condição; para (e), procurar um critério de pertinência.
\resolucao (a) Nenhum natural é negativo: vazio. (b) $x^2=9$ dá $x=3$ ou $x=-3$; com $x>0$, sobra $\{3\}$: unitário. (c) Não há inteiro entre 3 e 4: vazio. (d) Há infinitos reais entre 3 e 4 ($3{,}1$; $3{,}01$; $3{,}001$; \ldots): infinito. (e) Não há como decidir se $1000$ é \enquote{grande}: não define conjunto.
\comentario Compare (c) e (d): a mesma condição, vazio num universo e infinito no outro. Em (e), o problema não é o tamanho, e sim a falta de critério (erro 1-E4).
\end{exemplo}

\begin{exemplo}{}{3}{autoral}{1-H3}
\label{ex:u1:barbeiro}
Numa vila, há um barbeiro, homem e morador da vila, que barbeia exatamente os homens da vila que não se barbeiam a si mesmos. O barbeiro se barbeia?
\plano Testar as duas possibilidades contra a regra que descreve o barbeiro.
\resolucao Se o barbeiro se barbeia, ele é um homem que se barbeia; pela regra, o barbeiro \emph{não} o barbeia — isto é, ele não se barbeia. Contradição. Se ele não se barbeia, é um dos homens que não se barbeiam; pela regra, o barbeiro o barbeia — isto é, ele se barbeia. Contradição de novo. As duas possibilidades são impossíveis: \textbf{não existe} um homem da vila com essa propriedade.
\comentario Repare no que \emph{não} é contraditório: o conjunto $\{x\in\text{Homens da vila}\mid x\text{ não se barbeia}\}$ existe e pode ser descrito sem problema. O que não existe é um morador que barbeie exatamente esse conjunto. É a mesma saída de Zermelo para o paradoxo de Russell: separar elementos dentro de um universo dado nunca é contraditório; o argumento só mostra que certo objeto não está no universo.
\end{exemplo}

% ------------------------------------------------------------
\Questoes

\begin{questao}{1}{autoral}{1-H1}\label{u1:pertence}
Seja $A=\{1,\{2\},3\}$. Classifique em verdadeira ou falsa:
\begin{multicols}{4}
\begin{enumerate}[label=\alph*)]
  \item $2\in A$
  \item $\{2\}\in A$
  \item $3\in A$
  \item $\{3\}\in A$
\end{enumerate}
\end{multicols}
\noindent Quantos elementos tem $A$?
\resp{a) F;\ b) V;\ c) V;\ d) F. $A$ tem 3 elementos.}
\end{questao}

\begin{questao}{1}{autoral}{1-H4, 1-E3}\label{u1:quantos}
Quantos elementos tem cada conjunto?
\begin{multicols}{5}
\begin{enumerate}[label=\alph*)]
  \item $\varnothing$
  \item $\{0\}$
  \item $\{\varnothing\}$
  \item $\{\varnothing,\{\varnothing\}\}$
  \item $\{0,\varnothing\}$
\end{enumerate}
\end{multicols}
\resp{a) 0;\ b) 1;\ c) 1;\ d) 2;\ e) 2.}
\end{questao}

\begin{questao}{1}{autoral}{1-H2, 1-U1, 1-U2}\label{u1:listar}
Dê os elementos de cada conjunto:
\begin{enumerate}[label=\alph*)]
  \item os divisores naturais de 12;
  \item $\{x\in\mathbb{Z}\mid -2\le x<3\}$;
  \item os números primos menores que 20;
  \item as letras da palavra \emph{banana};
  \item o conjunto produzido pelo código Python \texttt{\{x for x in range(1, 20) if x \% 3 == 0\}}, sabendo que \texttt{range(1, 20)} percorre os inteiros de 1 a 19 e que \texttt{x \% 3 == 0} testa se $x$ é múltiplo de 3.
\end{enumerate}
\resp{a) $\{1,2,3,4,6,12\}$;\ b) $\{-2,-1,0,1,2\}$;\ c) $\{2,3,5,7,11,13,17,19\}$;\ d) $\{b,a,n\}$;\ e) $\{3,6,9,12,15,18\}$.}
\end{questao}

\begin{questao}{1}{adaptado do FME, exercício 13}{1-H2, 1-U1}\label{u1:fme13}
Escreva com símbolos:
\begin{enumerate}[label=\alph*)]
  \item o conjunto dos múltiplos inteiros de 3, entre $-10$ e $+10$;
  \item o conjunto dos divisores inteiros de 42;
  \item o conjunto dos múltiplos inteiros de 0;
  \item o conjunto dos números que se escrevem como uma fração $\frac{n}{d}$ com $n,d\in\{1,2\}$;
  \item o conjunto dos nomes das capitais dos estados da região Centro-Oeste do Brasil.
\end{enumerate}
\resp{a) $\{-9,-6,-3,0,3,6,9\}$;\ b) $\{\pm1,\pm2,\pm3,\pm6,\pm7,\pm14,\pm21,\pm42\}$;\ c) $\{0\}$;\ d) $\{\frac12,1,2\}$ (as frações $\frac11$ e $\frac22$ são o mesmo número);\ e) $\{$Cuiabá, Campo Grande, Goiânia$\}$.}
\end{questao}

\begin{questao}{1}{adaptado do FME, exercício 14}{1-H2, 1-U1}\label{u1:fme14}
Descreva cada conjunto por meio de uma propriedade dos seus elementos:
\begin{multicols}{2}
\begin{enumerate}[label=\alph*)]
  \item $A=\{+1,-1,+2,-2,+3,-3,+6,-6\}$
  \item $B=\{0,-10,-20,-30,-40,\ldots\}$
  \item $C=\{1,4,9,16,25,36,\ldots\}$
  \item $D=\{\text{Lua}\}$
\end{enumerate}
\end{multicols}
\resp{a) $\{x\in\mathbb{Z}\mid x$ é divisor de $6\}$;\ b) $\{x\in\mathbb{Z}\mid x$ é múltiplo de 10 e $x\le0\}$;\ c) $\{n^2\mid n\in\mathbb{Z},\ n>0\}$;\ d) $\{x\mid x$ é satélite natural da Terra$\}$ (universo: os corpos celestes). (Outras propriedades corretas também valem.)}
\end{questao}

\begin{questao}{2}{adaptado do FME, exercício 16}{1-H4, 1-U3, 1-E3}\label{u1:vazios}
No universo $\mathbb{Z}$, dizemos que um inteiro $a$ é \emph{divisível} por um inteiro $b$ (e que $b$ é \emph{divisor} de $a$) quando existe um inteiro $k$ tal que $a=b\cdot k$. Quais dos conjuntos abaixo são vazios?
\begin{multicols}{2}
\begin{enumerate}[label=\alph*)]
  \item $A=\{x\mid 0\cdot x=0\}$
  \item $B=\{x\mid x>\frac94\text{ e }x<\frac65\}$
  \item $C=\{x\mid x\text{ é divisor de zero}\}$
  \item $D=\{x\mid x\text{ é divisível por zero}\}$
\end{enumerate}
\end{multicols}
\resp{Só $B$. ($A=\mathbb{Z}$; $C=\mathbb{Z}$, pois $0=x\cdot0$ para todo $x$; $D=\{0\}$.)}
\end{questao}

\begin{questao}{2}{adaptado do FME, exercício 15}{1-H3, 1-U3, 1-E1}\label{u1:unitarios}
\begin{enumerate}[label=\alph*)]
  \item No universo $\mathbb{R}$, quais dos conjuntos abaixo são unitários?
\[
A=\Bigl\{x\ \Big|\ x<\tfrac94\text{ e }x>\tfrac65\Bigr\}\qquad B=\{x\mid 0\cdot x=2\}
\]
\[
C=\{x\mid x\text{ é inteiro e }x^2=3\}\qquad D=\{x\mid 2x+1=7\}
\]
  \item E se o universo for $\mathbb{Z}$?
\end{enumerate}
\resp{a) Só $D=\{3\}$.\ b) $A=\{2\}$ e $D=\{3\}$.}
\end{questao}

\begin{questao}[integradora]{2}{autoral}{1-H3, 1-U2, 1-E1}\label{u1:equacao}
Determine o conjunto das soluções de $2x^2-3x-2=0$ em $\mathbb{N}$, em $\mathbb{Z}$ e em $\mathbb{Q}$. Depois, determine o conjunto das soluções de $x^2+1=0$ em $\mathbb{R}$.
\resp{$\{2\}$ em $\mathbb{N}$;\ $\{2\}$ em $\mathbb{Z}$;\ $\{-\frac12,2\}$ em $\mathbb{Q}$;\ $\varnothing$ para $x^2+1=0$ em $\mathbb{R}$.}
\end{questao}

\begin{questao}{2}{autoral}{1-H3, 1-U2}\label{u1:distancia}
Seja $P$ um ponto. Descreva o conjunto dos pontos que estão à distância 1 de $P$ quando o universo é:
\begin{enumerate}[label=\alph*)]
  \item uma reta que passa por $P$;
  \item um plano que contém $P$;
  \item o espaço.
\end{enumerate}
\resp{a) dois pontos, um de cada lado de $P$;\ b) a circunferência de centro $P$ e raio 1;\ c) a superfície esférica de centro $P$ e raio 1.}
\end{questao}

\begin{questao}{2}{autoral}{1-H1, 1-E4}\label{u1:criterio}
Quais das descrições abaixo definem um conjunto? Para as que não definem, diga por quê.
\begin{enumerate}[label=\alph*)]
  \item os números naturais pequenos;
  \item os números naturais menores que $10^{100}$;
  \item os alunos da turma com média maior ou igual a 7;
  \item os bons alunos da turma;
  \item os estados brasileiros que fazem fronteira com o Uruguai.
\end{enumerate}
\resp{Definem conjunto: b, c, e (em e, o conjunto é $\{$Rio Grande do Sul$\}$). Não definem: a, d — \enquote{pequeno} e \enquote{bom} não decidem quem pertence.}
\end{questao}

\begin{questao}{2}{autoral}{1-H2, 1-U1, 1-E2}\label{u1:reticencias}
A lista $\{2,3,5,\ldots\}$ admite pelo menos duas leis de formação diferentes. Escreva duas, cada uma por meio de uma propriedade, e diga, em cada caso, se 8 pertence ao conjunto.
\resp{Por exemplo: $\{x\in\mathbb{N}\mid x$ é primo$\}$, ao qual 8 não pertence (depois do 5 vem o 7); e o conjunto dos termos da sequência $2,3,5,8,13,\ldots$, em que os dois primeiros termos são 2 e 3 e cada termo seguinte é a soma dos dois anteriores, ao qual 8 pertence ($3+5=8$).}
\end{questao}

\begin{questao}[integradora]{3}{autoral}{1-H3, 1-H4, 1-U3}\label{u1:parametro}
Considere, para cada número real $k$, o conjunto $A_k=\{x\in\mathbb{Z}\mid x^2-2x+k=0\}$. Para que valores de $k$ o conjunto $A_k$ é unitário? Para que valores tem exatamente dois elementos? E quando é vazio?
\resp{Unitário só para $k=1$ ($A_1=\{1\}$). Dois elementos quando $k=1-m^2$ com $m$ inteiro positivo ($k=0,-3,-8,-15,\ldots$; então $A_k=\{1-m,1+m\}$). Vazio para todos os outros valores de $k$.}
\end{questao}

\begin{questao}{3}{autoral}{1-H3, 1-U2}\label{u1:escolherU}
Seja $A=\{x\in U\mid (3x-2)(x^2-5)(x+4)=0\}$. Entre os universos $\mathbb{N}$, $\mathbb{Z}$, $\mathbb{Q}$ e $\mathbb{R}$, escolha um para o qual $A$ seja vazio; um para o qual $A$ seja unitário; e um para o qual $A$ tenha exatamente três elementos — se existir.
\resp{Vazio: $\mathbb{N}$. Unitário: $\mathbb{Z}$ ($A=\{-4\}$). Exatamente três elementos: nenhum dos quatro (em $\mathbb{Q}$, $A=\{-4,\frac23\}$ tem dois; em $\mathbb{R}$, $A=\{-4,\frac23,-\sqrt5,\sqrt5\}$ tem quatro).}
\end{questao}

\begin{questao}{3}{autoral}{1-H1, 1-E4}\label{u1:berry}
Considere a frase: \enquote{\emph{seja $n$ o menor número natural que não pode ser descrito com menos de vinte palavras}}. Mostre que essa frase não define um número.
\resp{Demonstração. A prova deve mostrar que, se a frase definisse um número $n$, ela própria — que tem menos de vinte palavras — o descreveria com menos de vinte palavras, contrariando a definição de $n$; logo a \enquote{propriedade} não tem critério de pertinência fixo.}
\end{questao}

\begin{sintese}
\textbf{Ao fim desta unidade, você deveria saber:}
\begin{itemize}[leftmargin=1.4em,itemsep=1pt]
  \item ler e escrever $x\in A$ e $x\notin A$, inclusive quando os elementos são conjuntos;
  \item passar de uma lista para uma propriedade e de uma propriedade para uma lista;
  \item fixar o universo antes de descrever um conjunto por propriedade;
  \item distinguir $\varnothing$, $\{0\}$ e $\{\varnothing\}$ e reconhecer conjuntos vazios e unitários;
  \item recusar descrições que não decidem quem pertence.
\end{itemize}
\smallskip\noindent\textbf{Responda com suas palavras, sem consultar.} Depois, confira com a caixa \emph{Para lembrar} (p.~\pageref{lembrar-u1}) e com a seção indicada.
\begin{enumerate}[label=\arabic*.,leftmargin=1.4em]
  \perguntasintese{O que é um conjunto, se ele não tem definição? \ondever{construção 1-H1}}
  \perguntasintese{Por que a descrição por propriedade precisa de um universo? \ondever{De onde vem?}}
  \perguntasintese{Como você reconhece, num enunciado, que precisa fixar o universo antes de resolver? \ondever{Onde é usado?}}
  \perguntasintese{Qual é o erro mais perigoso ao descrever um conjunto? Como você o detecta? \ondever{Erros comuns}}
\end{enumerate}
\end{sintese}

% ------------------------------------------------------------
\Respostas
{\small\resposta{S}{Síntese de domínio: confira com a caixa \emph{Para lembrar} (p.~\pageref{lembrar-u1}) e com a seção indicada em cada pergunta.}}
